%------------------------------------------------------------------------------
% \PART 1
%------------------------------------------------------------------------------
\chapter[Обзор существующих нейронных сетей для распознавания рукописных цифр]
{ОБЗОР СУЩЕСТВУЮЩИХ НЕЙРОННЫХ СЕТЕЙ ДЛЯ РАСПОЗНАВАНИЯ РУКОПИСНЫХ ЦИФР}

%------------------------------------------------------------------------------
% \PART 1.1 Text
%------------------------------------------------------------------------------
\section{Нейронные сети}\par
\hspace*{12.5 mm}Нейронная сеть представляет из себя совокупность нейронов,
соединенных друг с другом определенным образом. Нейрон представляет из себя 
элемент, который вычисляет выходной сигнал из совокупности входных сигналов.
Структура нейрона представлена на рисунке~\ref{fig:neiron}.

\insertfigure{neiron}{pic/neiron.png}{Общее представление нейрона}{10.0cm}

Нейроны выполняют следующую последовательность действий:

– приём входных сигналов;

– комбинирование входных сигналов;

– вычисление выходных сигналов;

– передача выходных сигналов следующим нейронам.

Между собой нейроны могут быть соединены абсолютно по-разному, это
определяется структурой конкретной сети. Искусственные нейронные сети 
можно рассматривать как направленный граф со взвешенными связями, в котором 
искусственные нейроны являются узлами. По архитектуре связей они могут быть 
сгруппированы в два класса: сети прямого распространения, в которых графы не 
имеют петель, и рекуррентные сети (RNN), или сети с обратными связями.

%------------------------------------------------------------------------------
% \PART 1.2 Text
%------------------------------------------------------------------------------
\section{Существующие нейронные сети для распознавания рукописных цифр}\par
\hspace*{12.5 mm}Распознавание рукописных цифр является одной из классических 
задач машинного обучения, для решения которой применяется широкий спектр 
нейронных сетей. В зависимости от сложности задачи, требований к точности, 
скорости обработки и аппаратных ограничений используются различные архитектуры, 
включая полносвязные сети (FCNN), сверточные нейронные сети (CNN) и более 
сложные гибридные модели.

Одним из первых методов распознавания рукописных цифр стали многослойные 
перцептроны (MLP), относящиеся к классу полносвязных нейронных сетей. 
Такие сети состоят из входного слоя, нескольких скрытых слоев и выходного слоя, 
где каждый нейрон соединен со всеми нейронами предыдущего и последующего слоев. 
Они способны успешно решать задачу распознавания, но 
требуют большого количества параметров и вычислительных ресурсов, что делает 
их менее эффективными.

Для обработки изображений, включая рукописные цифры, более эффективными 
оказались сверточные нейронные сети. Эти сети включают сверточные слои, 
которые извлекают характерные признаки из изображения. CNN значительно 
превосходят полносвязные сети в точности и скорости работы, поскольку 
используют пространственные зависимости в данных и требуют меньше параметров.

Хотя рекуррентные нейронные сети и их усовершенствованные версии, такие 
как Long Short-Term Memory Network (LSTM) и Gated Recurrent Unit Network (GRU),
чаще применяются для обработки последовательных данных, они могут 
использоваться и для распознавания рукописных цифр. В частности, они 
эффективны в случаях, когда рукописные символы анализируются в контексте 
строки.

Современные архитектуры, такие как Vision Transformer (ViT) и его модификации, 
предлагают альтернативный подход к обработке изображений, основанный на 
механизме самовнимания. Хотя традиционно трансформеры использовались в 
обработке естественного языка (NLP), их успешное применение в компьютерном 
зрении позволило достичь новых высот в распознавании символов и цифр.

Для повышения точности и эффективности могут применяться гибридные подходы, 
комбинирующие преимущества разных типов нейронных сетей. Например, модели, 
совмещающие CNN и LSTM, успешно применяются для распознавания рукописного 
текста, где CNN используется для выделения признаков, а LSTM – для анализа 
последовательностей.

В последующих разделах будет представлен детальный анализ наиболее 
распространённых нейросетевых моделей, применяемых для данной задачи, а также 
рассмотрены их основные особенности. Кроме того, будет приведена общая 
информация о различных функциях активации, используемых в современных 
архитектурах нейронных сетей.


%------------------------------------------------------------------------------
% \PART 1.3 Text
%------------------------------------------------------------------------------
\section{Полносвязные нейронные сети}\par
\hspace*{12.5 mm}Полносвязная нейронная сеть (Fully Connected Neural Network) — 
это базовая архитектура искусственных нейронных сетей, в которой каждый нейрон 
одного слоя соединен со всеми нейронами следующего слоя. Такая архитектура
чаще всего применяется в качестве классификатора после сверточных или 
рекуррентных слоев. Количество нейронов на каждом слое может отличаться от 
соседних слоев.

%------------------------------------------------------------------------------
% \PART 1.3.1 Text
%------------------------------------------------------------------------------
\anonsection{}Полносвязная нейронная сеть состоит из следующих основных 
компонентов:

    1 \text{Входной слой} – принимает данные.

    2 \text{Скрытые слои} – выполняют нелинейные преобразования входных данных. 
    Обычно включают несколько таких слоев.

    3 \text{Выходной слой} – содержит количество нейронов, соответствующее 
    числу классов, и использует функцию активации, 
    например softmax. Функция активации – это математическая функция, которая 
    определяет, передаст ли нейрон сигнал дальше по сети.

Математически данную сеть можно описать формулой для вычисления выхода нейрона 
в полносвязном слое
\begin{equation}
    {h} = f({W} {x} + {b}),
\end{equation}

\noindentгде   ${x}$ – входной вектор;\\
\hphantom{где} ${W}$ – матрица весов;\\
\hphantom{где} ${b}$ – вектор смещений;\\
\hphantom{где} $f(\cdot)$ – функция активации;\\
\hphantom{где} ${h}$ – выходной вектор нейронов.

Общая структура полносвязных нейронных сетей показана на 
рисунке~\ref{fig:fcnn}.

\insertfigure{fcnn}{pic/fcnn.png}{Общая структура полносвязных нейронных сетей}{14.0cm}

Для многослойной архитектуры выход слоя $l$ является входом для следующего слоя. 
В следствии чего, уравнения (1.1) можно преобразовать, с учетом многослойности, 
к виду
\begin{equation}
    {h}^{(l+1)} = f({W}^{(l)} {h}^{(l)} + {b}^{(l)}).
\end{equation}

\noindentгде   ${W^{(l)}}$ – матрица весов текущего слоя;\\
\hphantom{где} ${b^{(l)}}$ – вектор смещений текущего слоя;\\
\hphantom{где} ${h^{(l)}}$ – выходной вектор нейронов текущего слоя;\\
\hphantom{где} $f(\cdot)$ – функция активации;\\
\hphantom{где} ${h^{(l+1)}}$ – выходной вектор нейронов следующего слоя.

%------------------------------------------------------------------------------
% \PART 1.3.2 Text
%------------------------------------------------------------------------------
\anonsection{}ReLU (Rectified Linear Unit) – одна из наиболее популярных и 
широко используемых функций активации в нейронных сетях. Она определяется 
следующим образом

\begin{equation}
    f(x) = \max(0, x).
\end{equation}

\noindentгде   ${x}$ – входной вектор;\\
\hphantom{где} $max(\cdot)$ – операция выбора максимального значения из всех элементов в окне.

График функции ReLU представлен на рисунке~\ref{fig:relu}.

\insertfigure{relu}{pic/ReLU.png}{График функции ReLU}{8.5cm}

Функция ReLU ускоряет обучение сети и устранят проблему исчезающего градиента, 
так как имеет линейное поведение для положительных значений. Сама функция 
проста в реализации, но может быть чувствительна к большим значениям градиентов
и в отрицательной области возвращает 0.

%------------------------------------------------------------------------------
% \PART 1.3.3 Text
%------------------------------------------------------------------------------
\anonsection{}Сигмоида (Logistic Function) — это одна из классических 
функций активации, используемая в нейронных сетях. Математически представление 
в виде формулы

\begin{equation}
    f(x) = \frac{1}{1 + e^{-x}}.
\end{equation}

\noindentгде ${x}$ – входной вектор.

Сигмоида подходит для моделирования вероятностей, так как имеет диапазон 
значений $ (0,1) $. Функция гладкая и дифференцируемая, что упрощает вычисление
градиента, однако выходы не центрированы вокруг нуля, что замедляет обучение.

График функции сигмоида представлен на рисунке~\ref{fig:sigmoid}.

\insertfigure{sigmoid}{pic/sigmoid.png}{График функции сигмоида}{9cm}

%------------------------------------------------------------------------------
% \PART 1.3.4 Text
%------------------------------------------------------------------------------
\anonsection{}Функция Softmax применяется в многоклассовой классификации и 
преобразует входные значения в вероятностное распределение
\begin{equation}
    \sigma{(z_i)} = \frac{e^{z_i}}{\sum_{j} e^{z_j}},
\end{equation}

\noindentгде $z_i$ – входное значение для $i$-го класса.

Гистограмма функции Softmax представлен на рисунке~\ref{fig:sm}.

\insertfigure{sm}{pic/softmax.png}{Гистограмма функции Softmax}{8.0cm}

Softmax позволяет интерпретировать выходные значения сети как вероятности 
принадлежности к классам. Однако склонна к затуханию градиента при слишком 
больших входных значениях.

%------------------------------------------------------------------------------
% \PART 1.3.5 Text
%------------------------------------------------------------------------------
\anonsection{}Гиперболический тангенс (Tanh) — это сигмоидная функция, но 
симметричная относительно начала координат. Он определяется следующим образом

\begin{equation}
    f(x) = \tanh(x) = \frac{e^x - e^{-x}}{e^x + e^{-x}}.
\end{equation}

\noindentгде ${x}$ – входной вектор.

Функция принимает значения в диапазоне $(-1, 1)$, что делает её более 
предпочтительной по сравнению с сигмоидой в скрытых слоях нейронных сетей. 
Основное преимущество функции $\tanh(x)$ заключается в том, что производная 
функции принимает большие значения в среднем диапазоне, что уменьшает проблему 
исчезающего градиента по сравнению с сигмоидой.

Кроме того, благодаря симметрии относительно нуля, выходы функции $\tanh(x)$ 
сбалансированы — среднее значение ближе к нулю, что может способствовать более 
быстрой сходимости модели во время обучения. В отличие от сигмоидальной функции, 
где все значения положительные, $\tanh$ помогает избежать смещения активаций в 
одну сторону. Однако при больших по модулю входных значениях функция также 
подвержена эффекту насыщения, где градиенты становятся близкими к нулю.

График функции гиперболический тангенс представлен на рисунке~\ref{fig:tang}.

\insertfigure{tang}{pic/tang.png}{График функции гиперболический тангенс}{9.0cm}

%------------------------------------------------------------------------------
% \PART 1.3.6 Text
%------------------------------------------------------------------------------
\anonsection{}Обучение FCNN осуществляется с использованием метода обратного 
распространения ошибки (Backpropagation) и оптимизационного алгоритма, 
например, стохастического градиентного спуска (SGD)

\begin{equation}
    W^{(l)} \leftarrow W^{(l)} - \eta \frac{\partial L}{\partial W^{(l)}},
\end{equation}

\begin{equation}
    b^{(l)} \leftarrow b^{(l)} - \eta \frac{\partial L}{\partial b^{(l)}},
\end{equation}

\noindentгде   $\eta$ – скорость обучения;\\
\hphantom{где} $L$ – функция потерь, например, кросс-энтропия для классификаци.

Во время обратного распространения градиенты функции потерь вычисляются по 
отношению к весам и смещениям сети, после чего параметры обновляются в сторону, 
уменьшающую значение ошибки. Этот процесс повторяется итеративно для каждого 
батча данных, что позволяет модели постепенно улучшать свою способность к 
предсказанию.

Энтропия измеряет степень неопределенности в распределении вероятностей. 
Кросс-энтропия, в свою очередь, оценивает расхождение между двумя 
распределениями: истинным и предсказанным. Чем ближе $\hat{y}_i$ к $y_i$, тем 
меньше значение функции потерь.

Функцию кросс-энтропии описывается формулой 
\begin{equation}
    L = - \sum_{i} y_i \log(\hat{y}_i),
\end{equation}

\noindentгде   \( y_i \) – истинное значение (обычно one-hot вектор);\\
\hphantom{где} \( \hat{y_i} \) – вероятность, предсказанная моделью для класса 
$i$.

Кросс-энтропия особенно эффективна в задачах многоклассовой классификации, так как 
сильнее штрафует за уверенные, но неверные предсказания. Это способствует более 
быстрому обучению модели и лучшей сходимости.

%------------------------------------------------------------------------------
% \PART 1.3.7 Text
%------------------------------------------------------------------------------
\anonsection{}Структура однослойной нейронной сети прямого распространения, 
состоящей из полносвязного слоя с выходной функцией активации softmax показана 
на рисунке~\ref{fig:nn-struct}.

\insertfigure{nn-struct}{pic/NN_struct.png}{Однослойной нейронной сети прямого распространения}{12.5cm}

Данная архитектура позволяет добиться точности 92,4\%\cite{Doc}. Нейронная сеть 
использует 8624 параметра, которые необходимо хранить в памяти в качестве 
весовых коэффициентов и смещений.

%------------------------------------------------------------------------------
% \PART 1.4 Text
%------------------------------------------------------------------------------
\section{Сверточные нейронные сети}
\hspace*{12.5 mm}
Сверточные нейронные сети (Convolutional Neural Networks) представляют собой 
архитектуру глубокого обучения, предназначенную для обработки данных обладающих 
сетчатой топологией, таких как изображения. Их основное преимущество 
заключается в способности эффективно извлекать пространственные и иерархические 
особенности входных данных с помощью операций свертки.

Принцип работы сверточных нейронных сетей показан на рисунке~\ref{fig:cnn}\cite{CNN}.

\insertfigure{cnn}{pic/cnn.png}{Принцип работы сверточных нейронных сетей}{13cm}

%------------------------------------------------------------------------------
% \PART 1.4.1 Text
%------------------------------------------------------------------------------
\anonsection{}Стандартная архитектура CNN включает в себя следующие основные 
слои:

    1 \text{Сверточные слои} – выполняют операцию свертки, выделяя значимые 
признаки входных данных, такие как границы, текстуры и формы. 

    2 \text{Функция активации} – вводит нелинейность, позволяя сети 
аппроксимировать сложные зависимости.\@

    3 \text{Слои подвыборки (Pooling)} – уменьшают размерность представления, 
что снижает количество параметров и вычислительную сложность.

    4 \text{Полносвязные слои} – традиционные нейронные слои, выполняющие 
классификацию извлеченных признаков.

    5 \text{Функция потерь} – используется для оценки ошибки модели.

%------------------------------------------------------------------------------
% \PART 1.4.2 Text
%------------------------------------------------------------------------------
\anonsection{}Основная операция в сверточных нейронных сетях – это свертка, 
(convolution), которая заключается в пошаговом применении фильтра (ядра 
свёртки) к локальным областям входного изображения. Формально операция свёртки 
между входным изображением $X$ и ядром $K$ определяется следующим образом

\begin{equation}
    S(i, j) = \sum_{m} \sum_{n} X(i+m, j+n) \cdot K(m, n),
\end{equation}

\noindentгде  \( X(i, j) \) – входное изображение;\\
\hphantom{где } \( K(m, n) \) – ядро свертки;\\
\hphantom{где } \( S(i, j) \) – выходное изображение после свертки.

Сверточный слой применяется к входному изображению, используя несколько 
фильтров, каждый из которых извлекает определенные характеристики. Как правило,
чем глубже слой, тем более абстрактные признаки он способен захватывать.

%------------------------------------------------------------------------------
% \PART 1.4.3 Text
%------------------------------------------------------------------------------
\anonsection{}Для уменьшения размерности карт признаков применяется операция 
подвыборки. Один из наиболее распространенных методов – \text{max pooling}, 
который выбирает максимальное значение в окне \( k \times k \)
\begin{equation}
    P(i, j) = \max_{(m, n) \in k \times k} S(i+m, j+n).
\end{equation}

\noindent где \(P(i, j)\) – значение на выходе слоя подвыборки в позиции \((i, j)\);\\
\hphantom{где }\(S(i + m, j + n)\) – значение входной карты признаков в позиции \((i + m, j + n)\);\\
\hphantom{где }\(max(\cdot)\) – операция выбора максимального значения из всех элементов в окне.


Этот метод помогает уменьшить объем вычислений и делает модель более устойчивой
к сдвигам изображения.

%------------------------------------------------------------------------------
% \PART 1.4.4 Text
%------------------------------------------------------------------------------
\anonsection{}Функции потерь измеряют отклонение предсказаний модели от 
истинных значений. Выбор функции зависит от типа задачи:

    1 {Cross-Entropy Loss} — используется в задачах классификации, усиливает 
    влияние ошибочных предсказаний.

    2 {Mean Squared Error (MSE)} — применяется в регрессии, чувствительна к 
    выбросам.
    
    3 {Mean Absolute Error (MAE)} — альтернатива MSE, менее чувствительная к 
    выбросам.

    4 {Huber Loss} — сочетает MSE и MAE; более устойчив к выбросам.

    5 {Focal Loss} — модифицированная кросс-энтропия для задач с дисбалансом 
    классов.


%------------------------------------------------------------------------------
% \PART 1.4.5 Text
%------------------------------------------------------------------------------
\anonsection{}Обучение CNN происходит с использованием алгоритма обратного 
распространения ошибки и градиентного спуска. Для обновления параметров 
используется правило

\begin{equation}
    w^{(t+1)} = w^{(t)} - \eta \frac{\partial L}{\partial w},
\end{equation}

\noindentгде  \( \eta \) – коэффициент обучения;\\
\hphantom{где} \( \dfrac{\partial L}{\partial w} \) – градиент функции потерь 
по параметру \( w \).

%------------------------------------------------------------------------------
% \PART 1.4.6 Text
%------------------------------------------------------------------------------
\anonsection{}Структура сверточной нейронной сети, показатели точности 
которой 90\% показана на рисунке~\ref{fig:cnn-struct}\cite{Giardino}.

\insertfigure{cnn-struct}{pic/CNN_STRUCT.png}{Однослойной нейронной сети прямого распространения}{16cm}

Данная конфигурации состоит из: 

  1 Входной слой: [28$\times$28].

  2 Сверточный слой: 3 маски [3$\times$3]. 

  3 Слой пакетной нормализации. 

  4 Слой ReLU. 

  5 Полносвязанный слой. 

  6 Слой Softmax. 

  7 Слой классификации.

Нейронная сеть использует 23520 параметра, которые необходимо хранить в памяти 
в качестве весовых коэффициентов и смещений.

%------------------------------------------------------------------------------
% \PART 1.5 Text
%------------------------------------------------------------------------------
\section{Рекуррентные нейронные сети}
\hspace*{12.5 mm}Рекуррентные нейронные сети (Recurrent Neural Networks) 
представляют собой особый класс нейронных сетей, предназначенный для обработки 
последовательных данных, в которых важен порядок элементов во времени. 

В отличие от традиционных полносвязных или сверточных нейросетей, которые 
предполагают, что входы являются независимыми, RNN обладают внутренней памятью,
позволяющей учитывать контекст предыдущих состояний при формировании отклика на
текущий вход.

Принцип работы рекуррентной архитектуры нейронных сетей показан на 
рисунке~\ref{fig:rec}. Как видно из рисунка, ключевым элементом архитектуры 
является блок $Z^{-1}$, реализующий механизм обратной связи, за счёт 
которого выход текущего шага зависит не только от текущего входа.

\insertfigure{rec}{pic/rec.png}{Принцип работы рекуррентных нейронных сетей}{14cm}

%------------------------------------------------------------------------------
% \PART 1.5.1 Text
%------------------------------------------------------------------------------
\anonsection{}Основное отличие RNN от других типов нейронных сетей 
заключается в наличии обратных связей, позволяющих передавать информацию от 
предыдущих состояний. Для каждого временного шага \( t \) вычисляется новое 
состояние \( h_t \) на основе входных данных \( x_t \) и предыдущего состояния 
\( h_{t-1} \)
\begin{equation}
    h_t = f(W_x x_t + W_h h_{t-1} + b),
\end{equation}

\noindentгде \( h_t \) – скрытое состояние в момент времени \( t \);\\
\hphantom{где} \( x_t \) – входной вектор в момент времени \( t \);\\
\hphantom{где} \( W_x \), \( W_h \) – обучаемые весовые матрицы;\\
\hphantom{где} \( b \) – вектор смещения;\\
\hphantom{где} \( f \) – функция активации, например \( \tanh \) или \( ReLU \).

Выход сети \( y_t \) определяется как
\begin{equation}
    y_t = g(W_y h_t + c),
\end{equation}

\noindentгде \( W_y \) – матрица весов выходного слоя;\\ 
\hphantom{где} \( c \) – вектор смещения;\\ 
\hphantom{где} \( g \) – функция активации выходного слоя.

%------------------------------------------------------------------------------
% \PART 1.5.2 Text
%------------------------------------------------------------------------------
\anonsection{}Рекуррентные сети могут применяться для распознавания 
рукописных цифр. Один из распространенных подходов — обработка строк 
рукописного текста, где каждый символ представлен как последовательность 
пикселей или векторных признаков. В таких задачах используются LSTM (Long 
Short-Term Memory) или GRU (Gated Recurrent Unit), способные учитывать 
пространственно-временные зависимости между элементами изображения. 

Пример архитектуры LSTM представлен на рисунке~\ref{fig:rnn-struct}\cite{Varsamopoulos}.

\insertfigure{rnn-struct}{pic/rnn_struct.png}{Пример архитектуры LSTM}{9.5cm}

Данная архитектура позволяет решать задачу распознавания рукописного текста с 
точностью 94,7\%.

%------------------------------------------------------------------------------
% \PART 1.6 Text
%------------------------------------------------------------------------------
\section{Трансформеры нейронные сети}\par
\hspace*{12.5 mm}Трансформеры представляют собой архитектуру нейронных сетей, 
которые используют механизмы самовнимания (self-attention) для обработки 
входной информации параллельно, что позволяет достигать высокой эффективности и
точности.

%------------------------------------------------------------------------------
% \PART 1.6.1 Text
%------------------------------------------------------------------------------
\anonsection{}Архитектура трансформера состоит из энкодера и декодера, 
каждый из которых содержит несколько слоев. Основные компоненты трансформера:

  1 \text{Механизм самовнимания} позволяет модели учитывать 
важность различных частей входных данных.

  2 \text{Многоголовочное внимание (multi-head attention)} улучшает 
способность модели учитывать разные аспекты входной информации.

  3 \text{Обратная связь} ускоряет обучение и 
предотвращает исчезновение градиента.

  4 \text{Механизм нормализации} стабилизирует 
обучение.

  5 Полносвязные слои.

Механизм самовнимания вычисляется следующим образом
\begin{equation}
    \text{Attention}(Q, K, V) = \text{softmax}\left(\frac{QK^T}{\sqrt{d_k}}\right) V,
\end{equation}

\noindentгде $Q$, $K$ и $V$ — матрицы, полученные линейными 
преобразованиями данных;\\
\hphantom{где} $d_k$ — размерность ключей.

%------------------------------------------------------------------------------
% \PART 1.6.2 Text
%------------------------------------------------------------------------------
\anonsection{}Хотя трансформеры изначально были разработаны для задач 
обработки естественного языка, они также успешно применяются для распознавания 
рукописных цифр. Основной подход заключается в преобразование изображений в 
последовательности пикселей и обработка их с помощью Vision Transformer (ViT).

Трансформеры обеспечивают высокую точность  – 90,30\% и используют большое 
количество параметров – 163729 при обработке изображений, особенно при наличии 
больших объемов обучающих данных, что делает их перспективным направлением для 
распознавания рукописных цифр\cite{Agrawal}.

%------------------------------------------------------------------------------
% \PART 1.7 Text
%------------------------------------------------------------------------------
\section{Гибридные архитектуры}\par
\hspace*{12.5 mm}Гибридные архитектуры нейронных сетей представляют собой 
комбинацию различных типов нейронных сетей, объединяя их преимущества для 
улучшения производительности в задачах распознавания изображений, в том числе 
рукописных цифр. Такие модели могут включать элементы сверточных, рекуррентных, 
трансформерных и полносвязных сетей.

%------------------------------------------------------------------------------
% \PART 1.7.1 Text
%------------------------------------------------------------------------------
\anonsection{}Одним из наиболее распространенных гибридных подходов является 
использование сверточных нейронных сетей для извлечения признаков и 
рекуррентных нейронных сетей для обработки последовательностей. 

CNN выполняет обработку изображений, извлекая пространственные признаки, 
которые затем передаются в RNN.\@

RNN, такие как LSTM или GRU, анализируют временные зависимости между 
последовательными фрагментами извлеченных признаков, что особенно полезно 
при обработке последовательностей символов или цифр.

Математически процесс можно описать следующим образом
\begin{equation}
    F = \text{CNN}(X),
\end{equation}
\begin{equation}
    h_t = \text{RNN}(F_t, h_{t-1}),
\end{equation}

\noindentгде $X$ – входное изображение;\\ 
\hphantom{где} $F$ – извлеченные признаки;\\ 
\hphantom{где} $h_t$ – скрытое состояние RNN.\@

%------------------------------------------------------------------------------
% \PART 1.7.2 Text
%------------------------------------------------------------------------------
\anonsection{}Другой популярный гибридный подход – объединение CNN и 
трансформеров. CNN используются для локального извлечения признаков, а механизм
внимания трансформера помогает учитывать глобальный контекст изображения.

Формально, этот процесс можно записать следующим образом
\begin{equation}
    F = \text{CNN}(X),
\end{equation}
\begin{equation}
    Z = \text{Transformer}(F),
\end{equation}

\noindentгде $Z$ – закодированное представление входного изображения.

%------------------------------------------------------------------------------
% \PART 1.7.3 Text
%------------------------------------------------------------------------------
\anonsection{}Гибридные модели нейронных сетей находят широкое применение в 
задачах распознавания рукописных цифр благодаря своей способности эффективно 
обрабатывать как пространственные, так и временные зависимости. 

Такие 
архитектуры сочетают в себе преимущества различных типов слоёв – например, 
сверточных, рекуррентных и трансформеров, что позволяет достигать высокой 
точности классификации.

Одним из примеров является архитектура CNN-RNN, которая применяется для 
распознавания рукописных цифр, представленных в виде последовательностей, 
например, при обработке почтовых индексов. Сверточные слои извлекают 
пространственные признаки изображения, после чего рекуррентные слои 
анализируют последовательную структуру данных, позволяя учесть взаимосвязи 
между цифрами в строке.

Другой перспективный подход – использование CNN-Transformer моделей. Эти 
архитектуры применяются для высокоточной классификации изображений, в том 
числе зашумленных или искажённых. Благодаря высокой способности трансформеров к 
обобщению и контекстному анализу, удается достичь точности распознавания до 
99,89\% \cite{Performance}. 

Архитектура одной из таких моделей представлена на 
рисунке~\ref{fig:CNN-transformer}.

Благодаря комбинированию разных методов гибридные архитектуры обеспечивают 
улучшенные результаты в задачах компьютерного зрения.

\insertfigure{CNN-transformer}{pic/cnn_transformer.jpeg}{Архитектура CNN-Transformer}{16cm}

Входное изображение разбивается на небольшие фрагменты (patches). Это 
аналогично разбиению изображения на фиксированные блоки, которые затем 
используются как отдельные элементы входной последовательности.

Каждый патч преобразуется в вектор фиксированной размерности с помощью
линейного слоя. Также добавляется позиционное кодирование (Patch + Position 
Embedding), чтобы сохранить информацию о порядке патчей.

Вектора патчей подаются в трансформерный энкодер.

Затем он поступает в MLP Head, который выводит вероятности принадлежности 
к каждому из классов.